\section{Metodología}

La presente investigación se enmarca en un estudio de ..

\subsection{Síntesis y preparación de la ruta base (Ácido cítrico y urea)}
Para la ruta sintética base se empleó un sistema binario de ácido cítrico anhidro ($\text{C}_6\text{H}_8\text{O}_7$) y urea ($\text{CH}_4\text{N}_2\text{O}$). Para la relación molar 1:6, se pesaron $1{,}3910$~g de ácido cítrico y $2{,}6090$~g de urea, disolviéndose en 20~mL de agua desionizada bajo agitación constante a 350~rpm durante 15 minutos. La mezcla homogénea se sometió a pirólisis en un horno microondas a una potencia efectiva de 770~W durante 4 minutos, obteniéndose una matriz sólida monolítica (fase sólida). Para preparar la fase coloidal (líquida), una fracción del sólido se redisolvió en agua desionizada y se purificó mediante filtración por membrana de $0{,}22~\mu$m.
Para la relación molar 1:10, se mantuvo la masa de ácido cítrico, se incrementó proporcionalmente la cantidad de urea y se aplicó un tiempo de pirólisis de 8 minutos, obteniéndose una solución coloidal directamente luminosa.

\subsection{Síntesis de la ruta verde (Ácido cítrico y clara de huevo)}
Para la síntesis de la ruta verde sostenible se utilizó clara de huevo fresca (ovoalbúmina) como fuente orgánica de nitrógeno y azufre, combinada con ácido cítrico anhidro. La relación de partida consistió en $0{,}500$~g de ácido cítrico por cada $2{,}0$~mL de clara de huevo. Para garantizar la homogeneidad y reproducibilidad de toda la serie cinético-temporal, se preparó una solución madre multiplicada por un factor de escala $10\times$, disolviendo $5{,}000$~g de ácido cítrico anhidro en $20{,}0$~mL de clara de huevo y $20{,}0$~mL de agua desionizada. La solución madre se homogenizó bajo agitación magnética a 350~rpm durante 15 minutos.
Posteriormente, se tomaron alícuotas iguales de la solución madre y se sometieron a pirólisis por microondas a una potencia efectiva de 770~W, deteniendo la reacción a siete tiempos distintos: 1:00, 1:15, 1:30, 1:45, 2:00, 2:30 y 3:00~minutos. Tras la pirólisis, las alícuotas obtenidas se redisolvieron en agua desionizada y se purificaron por filtración mediante membrana de $0{,}22~\mu$m antes de la caracterización por espectroscopía UV-Vis.


\subsection{Fotometría digital RGB en cámara oscura}

Con el fin de obtener una caracterización complementaria de la emisión de
fotoluminiscencia de las muestras sintetizadas, se implementó un montaje
de fotometría digital RGB en cámara oscura. El montaje consistió en una
caja de cartón de dimensiones $30 \times 20 \times 12$~cm, cuyo interior
fue pintado con pintura acrílica negra mate para minimizar reflexiones
parásitas. Sobre la pared superior se fijó una linterna UV portátil con
longitud de onda de excitación $\lambda_{\text{exc}} = 365$~nm, cubriendo
la apertura con papel mantequilla como difusor para homogeneizar el haz
incidente sobre las muestras.

Las 11 muestras coloidales (Figura~\ref{fig:foto_viales}) se
distribuyeron dentro de la caja en tres filas horizontales de tubos
Eppendorf de 1.5~mL, dispuestos de la siguiente manera: fila superior
(V1 a V4), fila central (V5 a V8) y fila inferior (V9 a V11). La
adquisición de imágenes se realizó con un teléfono POCO X5 configurado
en modo manual con los siguientes parámetros: velocidad de obturación
$S = 1/2$~s, sensibilidad ISO~100, balance de blancos fijo en 5000~K y
lente ultra gran angular. Estos parámetros se mantuvieron constantes
durante toda la sesión de captura para garantizar la reproducibilidad
cuantitativa de los datos.

El procesamiento de las imágenes se llevó a cabo mediante un \textit{script}
en Python (OpenCV 5.0) que permitió la selección interactiva, por parte del
operador, del centro del líquido fluorescente de cada vial mediante clics
de mouse. Alrededor de cada centro seleccionado se definió una región de
interés (ROI) de tamaño fijo ($22 \times 44$~píxeles), garantizando que
el muestreo se realizara exclusivamente sobre el volumen líquido, sin
capturar las paredes del tubo ni el fondo oscuro. Para cada ROI se
extrajeron los valores promedio de los canales Rojo ($R$), Verde ($G$) y
Azul ($B$) del espacio de color sRGB. A partir de estos valores se
calcularon dos magnitudes:

\begin{enumerate}[label=\alph*)]
    \item La luminancia perceptiva según la recomendación BT.601:
    \begin{equation}
        Y = 0{,}299\,R + 0{,}587\,G + 0{,}114\,B
        \label{eq:luminancia}
    \end{equation}

    \item La relación verde-azul $G/B$, definida como el cociente entre
    los valores promedio del canal verde y el canal azul de cada ROI.
    Esta relación constituye un indicador adimensional del color dominante
    de la emisión, cuya ventaja fundamental es que resulta invariante
    frente a variaciones espaciales en la intensidad de la fuente de
    excitación, al tratarse de un cociente evaluado píxel a píxel dentro
    de la misma región.
\end{enumerate}